\documentclass[11pt,a4paper]{article}
\usepackage[utf8]{inputenc}
\usepackage[T1]{fontenc}
\usepackage{lmodern}
\usepackage[margin=2.4cm]{geometry}
\usepackage{amsmath,amssymb}
\usepackage{booktabs}
\usepackage{graphicx}
\usepackage{microtype}
\usepackage[ruled,vlined,linesnumbered]{algorithm2e}
\usepackage[round,authoryear]{natbib}
\usepackage[hidelinks]{hyperref}
\usepackage{enumitem}
\setlist{itemsep=2pt,topsep=4pt}

\title{\textbf{ExpertRuleFit: Deterministic Rule Ensembles\\with Protected Expert Rules for Regulated Decision-Making}}
\author{Djebar Hammouche\\Independent researcher, Luxembourg\\\texttt{github.com/hdjebar/ExpertRuleFit}}
\date{November 2026 --- technical report, preprint for review}

\begin{document}
\maketitle

\begin{abstract}
Rule ensembles such as RuleFit are readable, but their Lasso-based rule selection is unstable: changing the random seed changes the selected rules, which prevents an auditor from re-executing a model and obtaining the same output. Domain experts also need some rules to remain in the model even when the data signal is weak. We present \textsc{ExpertRuleFit}, a binary classifier that (i) generates candidate rules with a fixed-seed gradient-boosting ensemble, (ii) appends expert rules with reduced penalty weights, (iii) filters rules by their selection frequency across a fixed set of bootstrap elastic-net logistic fits, and (iv) enforces the inclusion of ``confirmatory'' expert rules through a constrained refit. On three public credit-scoring datasets, \textsc{ExpertRuleFit} returns bitwise-identical rule sets and predictions in every run (three runs per dataset, in separate processes), whereas RuleFit returns a different rule set for each of 100 seeds, with test AUC within 0.005 of RuleFit's mean. We state clearly that this benchmark measures \emph{determinism}, which holds by construction, and not stability under data perturbation, for which we propose a protocol. We discuss limitations, including the absence of sign constraints on expert rules, and a design for applying the method to tax-audit prioritisation.
\end{abstract}

\section{Introduction}

RuleFit \citep{friedman2008} turns the nodes of a tree ensemble into binary rules and selects a sparse subset with an $\ell_1$ penalty. The fitted model is an intercept plus a weighted sum of a few rules and linear terms, so each prediction can be traced to the conditions that produced it.

The selection step is unstable. Candidate rules are strongly correlated, and among correlated predictors the Lasso keeps one more or less arbitrarily. A different seed for the tree ensemble or for the cross-validation folds therefore gives a different rule set on the same data. In our benchmark (Section~\ref{sec:exp}), standard RuleFit returned a different rule set for each of 100 seeds, with a mean pairwise Jaccard similarity of 0.04--0.06.

In regulated settings such as credit, insurance or tax administration, this blocks deployment for three reasons:
\begin{itemize}
  \item \textbf{Re-execution.} A supervisor who re-runs the model must obtain the same output, otherwise decisions based on it cannot be audited.
  \item \textbf{Explanation.} A rule given to an analyst as the reason for a score should not disappear at retraining for reasons unrelated to the data.
  \item \textbf{Expert knowledge.} Criteria required by regulation or by business experts can be removed by the penalty when their signal is weak or collinear with other rules.
\end{itemize}

\paragraph{Contributions.}
\begin{enumerate}
  \item An estimator combining RuleFit rule generation, bootstrap selection-frequency filtering and penalty-weighted expert rules, with a constrained refit that guarantees the inclusion of confirmatory rules (Section~\ref{sec:method}).
  \item A precise statement of the reproducibility guarantee as determinism within an execution envelope, and of what it does not cover (Section~\ref{sec:repro}).
  \item A benchmark on three credit datasets, together with an explicit analysis of what it does and does not show, and a protocol for measuring stability under data perturbation (Section~\ref{sec:exp}).
  \item An open-source Python implementation (about 4{,}000 lines, 39 automated tests) built on \texttt{imodels} \citep{singh2021} and \texttt{scikit-learn}.
\end{enumerate}

\section{Related work}\label{sec:related}

\paragraph{Stable rule ensembles.} SIRUS \citep{benard2021} obtains stability from the forest structure: splits are discretised on quantiles and the most frequent rules across trees are kept. Two runs on the same data yield nearly the same rule list, and stability under data perturbation is reported. Other methods target sparsity: Compressed Rule Ensembles \citep{nalenz2022} merge similar rules, FIRE \citep{liu2023} uses a fusion penalty favouring shared antecedents, LIRE \citep{kanamori2023} minimises the number of rules needed per prediction, and FIGS \citep{tan2025} grows a small sum of trees. The \texttt{pre} package \citep{fokkema2020} uses relaxed and adaptive Lasso to reduce false selections.

\paragraph{Stability selection.} Retaining variables selected in most subsample fits is the stability selection of \citet{meinshausen2010}, with improved error bounds in \citet{shah2013}. \textsc{ExpertRuleFit} applies the idea to rule ensembles with a fixed threshold and ten draws, without the error control these works provide.

\paragraph{Prior knowledge in penalties.} Reducing the penalty on chosen variables follows the adaptive Lasso \citep{zou2006} and the use of per-feature penalty factors. Monotonicity constraints, common in credit scoring, fix the direction of an effect; \textsc{ExpertRuleFit} enforces inclusion, not direction.

\paragraph{Glass-box additive models.} Explainable Boosting Machines \citep{lou2013,nori2019} fit additive shape functions with pairwise interactions; we use them to propose interaction rules and in a stacked variant.

\paragraph{Model multiplicity.} Many models fit the data almost equally well \citep{rudin2024}. In that view, RuleFit's instability is a random draw from the Rashomon set, and \textsc{ExpertRuleFit} fixes the draw. This makes the output auditable but does not show that the chosen model is better than its neighbours.

\begin{table}[h]
\centering\small
\caption{Positioning (summary of the cited works).}
\label{tab:position}
\resizebox{\textwidth}{!}{\begin{tabular}{@{}lllll@{}}
\toprule
Method & Rule ensemble & Stable under perturbation & Deterministic re-run & Expert rules protected\\
\midrule
RuleFit & yes & no & only if all seeds fixed & no\\
SIRUS & yes & yes (reported) & near-identical & no\\
CRE, FIRE, LIRE & yes & not the main goal & solver-dependent & no\\
Stability selection & no & yes, with bounds & with fixed subsamples & no\\
\textsc{ExpertRuleFit} & yes & not yet measured & yes, by design & inclusion, not sign\\
\bottomrule
\end{tabular}}
\end{table}

\section{Method}\label{sec:method}

Let $X\in\mathbb{R}^{n\times p}$ and $y\in\{0,1\}^n$. \textsc{ExpertRuleFit} fits a logistic model on an augmented matrix
\[
Z = [\,x_1,\dots,x_p,\; r_1(X),\dots,r_k(X),\; c_1(X),\dots,c_m(X),\; o_1(X),\dots,o_q(X)\,],
\]
where $r_j$ are tree-derived binary rules, $c_j$ confirmatory expert rules and $o_j$ optional expert rules. Predictions are $P(y=1\mid x)=\sigma(\beta_0+\sum_j \beta_j z_j)$.

\paragraph{Standardisation and penalty weighting.} Columns are first standardised with statistics from the training data: linear terms are centred and divided by their standard deviation $s_j$, rule columns are divided by $s_j$ only, which keeps them sparse. Each column then receives a weight $w_j$ (1 for linear terms and tree rules, $w_c=10^{-2}$ for confirmatory rules, $w_o=0.3$ for optional rules by default) and is rescaled as $\tilde z_j = (z_j-m_j)/(s_j\sqrt{w_j})$ before an elastic-net fit, which scales the effective penalties on the standardised coefficients:
\[
\text{L1 term}\propto \sqrt{w_j}\,|\beta_j|,\qquad \text{L2 term}\propto w_j\,\beta_j^2 .
\]
With $w_c=10^{-2}$ the $\ell_1$ penalty on confirmatory rules is reduced by a factor of 10 and the $\ell_2$ penalty by a factor of 100. Coefficients and the intercept are mapped back to the original scale before reporting. Earlier versions used $w_c=10^{-8}$ without standardisation. A confirmatory column was then $10^4$ times larger than the others, which made the saga solver stop far from the optimum without a convergence warning. After standardisation, the conditioning of the problem depends on the ratio $1/\sqrt{w_c}$ rather than on the data; $w_c=10^{-2}$ is the smallest value that fitted without convergence warnings in the sensitivity study below. Inclusion of confirmatory rules does not rely on $w_c$: they bypass the frequency filter, and the refit below restores any the solver sets to zero.

\begin{algorithm}[h]
\caption{\textsc{ExpertRuleFit} training}\label{alg:erf}
\KwIn{$X$, $y$, confirmatory rules $C$, optional rules $O$, $B=10$, threshold $\tau=0.8$}
Fit gradient boosting (250 trees, depth 4, seed 42) and extract at most 50 rules $R$ via \texttt{imodels} RuleFit\;
Build $Z$ from linear terms, $R$, $C$, $O$; set weights $w$\;
\For{$b=1$ \KwTo $B$}{
  Draw bootstrap sample with seed $42+b$\;
  Fit elastic-net \texttt{LogisticRegressionCV} on $Z/\sqrt{w}$ (saga, 3-fold, $\ell_1$ ratio in $\{0.1,\dots,1\}$)\;
  Record which coefficients are non-zero\;
}
Keep columns selected in at least $\tau B$ fits, plus all of $C$\;
Fit final elastic-net \texttt{LogisticRegressionCV} on kept columns\;
\If{some $c\in C$ has $\beta_c=0$}{Refit with $C$ unpenalised\;}
\KwOut{rules, coefficients, flag \texttt{confirmatory\_all\_active\_}}
\end{algorithm}

The user's \texttt{random\_state} is ignored (a warning is emitted); all internal seeds derive from a fixed base seed. Two extensions are provided: \texttt{discover\_interaction\_rules}, which uses an EBM to find interacting feature pairs and thresholds and returns them as \emph{optional} rules, and \texttt{DualModel}, which stacks EBM and \textsc{ExpertRuleFit} scores with a logistic meta-model trained on out-of-fold predictions. Optional post-hoc calibration (Platt or isotonic, 5-fold) is available but off by default.

\paragraph{Expert-imposed versus data-discovered rules.} Confirmatory status is reserved for rules an expert imposes. Rules found by the EBM extension are data-discovered: they enter as optional rules, with a reduced penalty, and must pass the frequency filter like any tree rule. An expert can promote one to confirmatory after review; each discovered rule records its provenance (\texttt{source="ebm\_discovered"}) for the audit trail.

\paragraph{Sign check and absorption diagnostics.} A confirmatory rule may declare an expected sign $s_c\in\{+1,-1\}$. After fitting, the estimator compares $\operatorname{sign}(\beta_c)$ with $s_c$, sets a flag \texttt{confirmatory\_sign\_ok\_} and emits a warning for each contradiction; the sign is checked, not enforced. For each confirmatory rule it also reports, on the training data, the largest absolute correlation with any other candidate column, its variance inflation factor among retained columns, the candidate columns with $|r|\ge 0.5$ that are inactive in the final model (removed by the frequency filter or set to zero), and the change in the coefficients of retained correlated columns when the confirmatory rule is removed (ridge refits with the final model's $C$ and the same weights). These diagnostics make absorption visible; they do not prevent it.

\section{Reproducibility guarantee}\label{sec:repro}

The guarantee is bitwise determinism within an execution envelope: the same training data, in the same column order, on the same software stack, yields the same rule set and predictions. The envelope requires pinned versions of Python, NumPy, scikit-learn and imodels; single-threaded BLAS and OpenMP; the fixed internal seeds; a validated input schema with stable feature order; and identical platform settings. A container recipe with pinned image digests is provided in the repository.

The guarantee does not cover different data (adding or correcting rows may change the rule set), a different software stack or CPU, or dependence on row order where underlying algorithms break ties by position. Determinism makes re-execution auditable; robustness to small data changes, which governs whether explanations survive retraining, is a distinct property.

\section{Experiments}\label{sec:exp}

\paragraph{Data and protocol.} We use German Credit (UCI, 1{,}000 rows), Taiwan Credit (UCI; a random subsample of 5{,}000 of the 30{,}000 rows, seed 42, to bound runtime) and HMEQ (5{,}960 rows). Each dataset is split once into 80\% training and 20\% test data (stratified, seed 42). Missing values are imputed (median, or mode for categorical columns), categorical columns are ordinally encoded and all columns are standardised. All preprocessing statistics are fitted on the training split. RuleFit (imodels, 250 trees of depth 4, at most 50 rules) is fitted 100 times with \texttt{random\_state} $=0,\dots,99$. \textsc{ExpertRuleFit} (benchmark settings, solver tolerance $10^{-6}$) is fitted three times with \texttt{random\_state} $=0,1,2$, each run in a separate process; since the seed is ignored by design, more runs would repeat the same computation, and we instead check that the three runs agree bit for bit, comparing a hash of the rule set and the SHA-256 digest of the test-set probabilities. A run's rule set is the set of terms with non-zero coefficient; we report how many RuleFit runs share the most frequent rule set and the Jaccard similarity averaged over all pairs of runs. All runs use the pinned software stack (Python 3.11.18, scikit-learn 1.7.2, NumPy 2.4.6, imodels 3.0.4, single-threaded), in the container published with the code. An earlier version of this table fitted the preprocessing on the full dataset before the split and used an unpinned stack.

\begin{table}[h]
\centering\small
\caption{Seed sensitivity and discrimination (single train/test split). RuleFit: 100 seeds; AUC is the mean over seeds. \textsc{ExpertRuleFit}: three runs with bitwise-identical output.}
\label{tab:results}
\begin{tabular}{@{}lcccc@{}}
\toprule
 & \multicolumn{2}{c}{Runs with the most frequent rule set} & \multicolumn{2}{c}{Test AUC}\\
\cmidrule(lr){2-3}\cmidrule(lr){4-5}
Dataset & RuleFit & \textsc{ExpertRuleFit} & RuleFit & \textsc{ExpertRuleFit}\\
\midrule
German Credit & 1/100 & 3/3 & 0.788 ($\pm$0.009) & 0.790\\
Taiwan Credit & 1/100 & 3/3 & 0.764 ($\pm$0.003) & 0.769\\
HMEQ & 1/100 & 3/3 & 0.903 ($\pm$0.003) & 0.903\\
\bottomrule
\end{tabular}
\end{table}

Mean pairwise Jaccard similarity of RuleFit rule sets is 0.045 (German Credit), 0.058 (Taiwan Credit), 0.036 (HMEQ); on each dataset the 100 RuleFit runs produced 100 distinct rule sets. The three \textsc{ExpertRuleFit} runs per dataset produced the same rule set (43, 42, 48 terms) and the same probabilities. One \textsc{ExpertRuleFit} fit took between about 10 minutes (German Credit) and 3.5 hours (HMEQ) single-threaded, against seconds for RuleFit.

\paragraph{What the results show.} RuleFit's output depends heavily on the seed, so a single run cannot be reproduced without recording every seed. \textsc{ExpertRuleFit} removes this dependence without loss of discrimination on these datasets.

\paragraph{What they do not show.} Because \textsc{ExpertRuleFit} ignores \texttt{random\_state}, varying it does not change the computation: identical output across runs follows from the design and is not an empirical finding about stability; the three runs only confirm that the implementation is deterministic within the execution envelope. The comparison sets determinism against seed sensitivity. AUC is also reported on a single split without confidence intervals.

\paragraph{Penalty-weight sensitivity (German Credit).} We fitted \textsc{ExpertRuleFit} on the same split for $w_c\in\{10^{-8},10^{-6},10^{-4},10^{-2},1\}$ and $w_o\in\{0.1,0.3,1\}$ on the pinned software stack. The expert rules are standard credit-risk criteria chosen by the author, not rules supplied by domain experts: a confirmatory rule ``checking account below 0~DM'' and an optional rule ``duration above 24 months'', both expected to raise the default probability. Saga raised 360--361 convergence warnings per fit at $w_c=10^{-8}$, 66--69 at $10^{-6}$, 3--6 at $10^{-4}$ and none at $10^{-2}$ or $1$. Test AUC was 0.753--0.756 at $w_c=10^{-8}$ and 0.784--0.797 for all larger values, and fits at $10^{-2}$ took about a third of the time of fits at $10^{-4}$. The confirmatory coefficient had the expected positive sign only at $w_c\le 10^{-6}$, in fits that had not converged; this agreement is an artefact of non-convergence. In every converged fit the coefficient was negative: the model also retains the ordinal checking-status variable and tree rules containing the same condition, which carry the effect, and the sign check flags the contradiction. The optional rule was retained in every fit. These results cover one dataset and one split.

\paragraph{Proposed perturbation protocol.} To measure stability under data perturbation we will: (1) draw 100 bootstrap or 90\% subsamples of the training set; (2) fit RuleFit with a fixed seed and \textsc{ExpertRuleFit} on each; (3) report mean pairwise Jaccard similarity and per-rule selection frequencies, as in \citet{benard2021}; (4) report AUC with confidence intervals over repeated splits; (5) include SIRUS and $\ell_2$-penalised logistic regression as baselines. These results are not yet available.

\section{Limitations}\label{sec:limits}

\begin{itemize}
  \item \textbf{Stability under perturbation is unmeasured} (Section~\ref{sec:exp}).
  \item \textbf{No sign constraint.} A confirmatory rule is guaranteed to stay in the model, not to keep the direction the expert intended. A declared expected sign is checked and a contradiction is flagged, but the coefficient is not constrained. On German Credit the contradiction occurred in every converged fit (Section~\ref{sec:exp}); a sign-constrained, unpenalised confirmatory block is the next planned change.
  \item \textbf{Weakly penalised rules can absorb signal.} With a small $w_c$, a confirmatory rule correlated with tree rules can take their weight; a weak or wrong expert rule then shapes the model. The diagnostics report this case; they do not correct it.
  \item \textbf{Calibration not evaluated.} Penalised logistic outputs are often too conservative; optional post-hoc calibration exists but reliability curves and Brier scores have not been reported.
  \item \textbf{Coarse frequency estimate.} With $B=10$, the 80\% threshold means 8 of 10 draws; the draws are fixed, not resampled.
  \item \textbf{Scope.} Binary classification only; evidence limited to three credit datasets of moderate size; rule readability (up to four conditions) not tested with users.
  \item \textbf{Review.} Single author, internal code review only; the implementation was written with AI assistance.
\end{itemize}

\section{Application design: tax-audit prioritisation}\label{sec:app}

A tax administration ranking returns for review needs a reason an agent can state, a selection that can be replayed if contested, and inspectors' criteria kept in the model. We sketch a design; the method has not been applied to tax data.

\begin{table}[h]
\centering\small
\caption{Mapping of model elements to a tax-audit setting (illustrative).}
\begin{tabular}{@{}p{4cm}p{10.5cm}@{}}
\toprule
Model element & Tax-audit counterpart\\
\midrule
Target $y$ & Audit outcome, e.g.\ adjustment above a threshold set with the business\\
Tree-derived rules & Patterns in past audit results, e.g.\ high deductions combined with mainly salaried income\\
Confirmatory rules & Legal or business criteria, e.g.\ a gap between declared salary and the employer's statement, or deductions above a statutory ceiling\\
Optional rules & Inspectors' hypotheses, kept only if supported by data\\
Determinism & The selection list can be replayed for a contested case\\
Active rules per case & The explanation given to the reviewing agent\\
\bottomrule
\end{tabular}
\end{table}

Three conditions precede any pilot. First, \emph{selection bias}: outcomes exist only for audited files, so a naive model learns past selection habits; a small random audit sample and inverse-probability weighting are needed. Second, \emph{fairness and legal grounds}: variables without legal basis are excluded and score distributions compared across taxpayer groups. Third, \emph{human decision}: the score suggests a file and an agent decides, consistent with Article 22 GDPR on automated decisions. A pilot would compare the hit rate per audit hour with the current selection method on the random sample.

\section{Conclusion and future work}

\textsc{ExpertRuleFit} makes rule-ensemble outputs exactly reproducible and keeps expert-imposed rules, at no loss of AUC on three credit datasets. Its central empirical claim so far is determinism, not stability under data change. Future work, in order of priority: the perturbation protocol of Section~\ref{sec:exp} with SIRUS as baseline; sign-constrained confirmatory rules via bounded optimisation; a principled choice of the confirmatory weight; more bootstrap draws or error-controlled thresholds \citep{shah2013}; calibration assessment; and handling of selection bias when labels come only from past audits.

\paragraph{Code and data availability.} Code, tests, the validation script and a container recipe are at \url{https://github.com/hdjebar/ExpertRuleFit}. Datasets are public.

\begin{thebibliography}{99}\small
\bibitem[B{\'e}nard et~al.(2021)]{benard2021} B{\'e}nard, C., Biau, G., da~Veiga, S. and Scornet, E. (2021). SIRUS: Stable and interpretable RUle set for classification. \emph{Electronic Journal of Statistics}, 15(1), 427--505.
\bibitem[Fokkema(2020)]{fokkema2020} Fokkema, M. (2020). Fitting prediction rule ensembles with R package pre. \emph{Journal of Statistical Software}, 92(12).
\bibitem[Friedman and Popescu(2008)]{friedman2008} Friedman, J.~H. and Popescu, B.~E. (2008). Predictive learning via rule ensembles. \emph{Annals of Applied Statistics}, 2(3), 916--954.
\bibitem[Kanamori(2023)]{kanamori2023} Kanamori, K. (2023). Learning locally interpretable rule ensemble. In \emph{ECML PKDD}.
\bibitem[Liu and Mazumder(2023)]{liu2023} Liu, B. and Mazumder, R. (2023). FIRE: An optimization approach for fast interpretable rule extraction. In \emph{KDD}.
\bibitem[Lou et~al.(2013)]{lou2013} Lou, Y., Caruana, R., Gehrke, J. and Hooker, G. (2013). Accurate intelligible models with pairwise interactions. In \emph{KDD}.
\bibitem[Meinshausen and B{\"u}hlmann(2010)]{meinshausen2010} Meinshausen, N. and B{\"u}hlmann, P. (2010). Stability selection. \emph{Journal of the Royal Statistical Society B}, 72(4), 417--473.
\bibitem[Nalenz and Augustin(2022)]{nalenz2022} Nalenz, M. and Augustin, T. (2022). Compressed rule ensemble. In \emph{AISTATS}, PMLR 151.
\bibitem[Nori et~al.(2019)]{nori2019} Nori, H., Jenkins, S., Koch, P. and Caruana, R. (2019). InterpretML: A unified framework for machine learning interpretability. arXiv:1909.09223.
\bibitem[Rudin et~al.(2024)]{rudin2024} Rudin, C. et~al. (2024). Amazing things come from having many good models. In \emph{ICML}.
\bibitem[Shah and Samworth(2013)]{shah2013} Shah, R.~D. and Samworth, R.~J. (2013). Variable selection with error control: another look at stability selection. \emph{Journal of the Royal Statistical Society B}, 75(1), 55--80.
\bibitem[Singh et~al.(2021)]{singh2021} Singh, C., Nasseri, K., Tan, Y.~S., Tang, T. and Yu, B. (2021). imodels: a Python package for fitting interpretable models. \emph{Journal of Open Source Software}, 6(61).
\bibitem[Tan et~al.(2025)]{tan2025} Tan, Y.~S., Singh, C., Nasseri, K., Agarwal, A. and Yu, B. (2025). Fast interpretable greedy-tree sums. \emph{PNAS}.
\bibitem[Zou(2006)]{zou2006} Zou, H. (2006). The adaptive lasso and its oracle properties. \emph{Journal of the American Statistical Association}, 101(476), 1418--1429.
\end{thebibliography}

\end{document}
